\documentclass[notheorems]{beamer}
\mode<presentation>
%------idioma-------
\usepackage[spanish,mexico]{babel}
\usepackage{ragged2e}
\uselanguage{spanish}
\languagepath{spanish}

%----mates-------
\usepackage{amsmath,amsfonts}
\usepackage{amsthm}
\usepackage{pifont}
\usepackage{amssymb,wasysym}
\usepackage{mathtools}
\usepackage{bibentry}
\usepackage{pgfgantt}
%-----Color fuente-----
\usepackage{xcolor,colortbl}
\definecolor{UADY_azul}{rgb}{0,0.265, 0.484}
\definecolor{UADY_dorado}{rgb}{0.83,0.62, 0.05}
%-----Nuevos comandos-----
\newcommand{\C}{\mathbb{C}}
\newcommand{\R}{\mathbb{R}}
\newcommand{\Q}{\mathbb{Q}}
\newcommand{\Z}{\mathbb{Z}}
\newcommand{\N}{\mathbb{N}}
\newcommand{\Sf}{\mathbb{S}}
\newcommand{\br}{\langle}
\newcommand{\ke}{\rangle}
%-----Cajitas y afines-----
\usepackage{float}
\usepackage[many]{tcolorbox}
\tcbuselibrary{breakable}
\tcbuselibrary{theorems}
\usepackage{geometry}
\newtheorem{teo}{Teorema}[section]
\newtheorem{cor}[teo]{Corolario}
\newtheorem{lema}[teo]{Lema}
\newtheorem{prop}[teo]{Proposición}
\newtheorem{defi}[teo]{Definición}
\newtheorem{obs}[teo]{Observación}
\newtheorem{ejem}[teo]{Ejemplo}
\newtheorem{ejer}[teo]{Ejercicio}
%-------TikZ----------
\usepackage{tikz}
\usepackage{pgfplots}
\pgfplotsset{compat=1.18}
\usepackage{mathdots}
\usepackage{yhmath}
\usepackage{cancel}
\usepackage{color}
\usepackage{multirow}
\usepackage{array, tabularx}
\usepackage{extarrows}
\usepackage{booktabs}
\usetikzlibrary{fadings}
\usetikzlibrary{patterns}
\usetikzlibrary{shadows.blur}
\usetikzlibrary{shapes}

%---Referencias---
\usepackage{csquotes}
\usepackage[
backend=biber,
style=alphabetic,
sorting=ynt
]{biblatex}
\addbibresource{Referencias.bib}
\renewcommand{\bibfont}{\normalfont\footnotesize}
\usepackage{hyperref}
\hypersetup{colorlinks=true,linkcolor=UADY_azul,urlcolor=UADY_azul,citecolor=UADY_azul}
\urlstyle{same}
%----tema de la presentación------
\usetheme{Berkeley}
\logo{\includegraphics[width=0.85 cm,keepaspectratio]{UADYLOGO.png}}
\usecolortheme[named=UADY_dorado]{structure}
\setbeamertemplate{blocks}[rounded][shadow=true]
\setbeamertemplate{items}[ball]
\setbeamercolor{bibliography item}{fg=UADY_azul}
\setbeamercolor{bibliography entry author}{fg=UADY_azul}
\setbeamercolor{bibliography entry title}{fg=UADY_azul}
\setbeamercolor{bibliography entry note}{fg=UADY_azul}
\renewcommand{\familydefault}{\sfdefault}
    \title{Título de la plática}
    
        \subtitle{Subtítulo de la plática}
    
    \author{Nombre(s) del(los) autor(es)}


    \institute{Universidad Autónoma de Yucatán}

    \date{\today}

\newcounter{currentenumi} 

\begin{document}

\frame{\titlepage}

\frame{\tableofcontents}



\section{Nombre de la sección}

\begin{frame}{Título de la diapositiva}
\begin{itemize}
    \item Aquí una ecuación:
\pause
\begin{equation*}
I_{a+}^{\alpha}f(x)=\frac{1}{\Gamma(\alpha)}\int_a^x(x-t)^{\alpha-1}f(t)dt.
\end{equation*}
\vspace{0.5 cm}
\pause
\item Aquí otra ecuación con símbolos previamente renombrados en el preámbulo; tu puedes renombrar otros a tu gusto:
\begin{equation*}   \R,\Z,\Q,\C,\Sf^n,\N,\br\cdot,\cdot\ke
\end{equation*}
\end{itemize}
\end{frame}
\section{Nombre de la sección}
\begin{frame}{Título de la diapositiva}
    \begin{itemize}
        \item Aquí hago referencia a \cite{stein2010complex}.
        \pause\item Ahora un teorema:
        \pause\begin{teo}\label{teo_ejemplo}
            Si $A\in \R^{n\times n} ~(n\in\N)$ y $p_A$ es el polinomio característico de $A$, entonces $p_A(A)=O$, siendo $O$ la matriz de ceros. 
        \end{teo}
    \end{itemize}
\end{frame}
\section{Nombre de la sección}
\begin{frame}{Titulo de la diapositiva}
    \begin{figure}
        \centering
        \input{Figura_ejemplo}
        \caption{Ejemplo de una figura.}
        \label{fig_ejemplo}
    \end{figure}
\end{frame}
\section{Referencias}
\begin{frame}
\printbibliography    
\end{frame}

\end{document}